\section{Remaining questions}\label{sec:open}
This section identifies the wild-index and rational-torsion questions left open by the proved results. Theorem~\ref{thm:complete-prime-index} determines every full-class height-two index with $p\nmid e$ and $\mu_p\subset K$. The optimal supersingular-only range beyond Theorem~\ref{thm:ss-iterated-range} is a separate question, since the sharp full-class lower examples can be mixed comparisons.

At indices divisible by $p$, Theorem~\ref{thm:first-wild-general} gives $\nuMin(p,p(p-1),2)=2$ for every $p\ge5$, and Theorem~\ref{thm:wild-bands-exact} determines every coefficient $2\le k\le(p^2-1)/2$ in $e=kp^r(p-1)$: the index is $r+2$ when $p\nmid k$ or $k=p$, and $r+3$ for the other multiples of $p$. Theorem~\ref{thm:all-index-upper} gives the uniform sufficient range $e\le p^{n-1}$ for every index, and Theorem~\ref{thm:mixed-all-indices} determines the exact mixed-comparison threshold at every index. The remaining full-class problem concerns comparisons between two whole supersingular groups outside these proved ranges. More precisely, Corollary~\ref{cor:all-index-two-slope-reduction} reduces recovery under
\[
 e\le p^{2n-3-v_p(e)}
\]
to two models with two first-level Newton slopes. Theorem~\ref{thm:all-index-central-reduction} further produces a central comparison at a remaining level $p^r$, where $r\ge v_p(e)+2$ and the high numerator has exact $p$-adic valuation $v_p(e)$. The partner polygon is unrestricted, and the original numerical bound is not automatically valid at level $r$. For $v_p(e)>0$, exclusion of these terminal configurations remains open. Theorem~\ref{thm:full-supersingular-maximum} expresses the full index exactly as the maximum of the supersingular-only index and the mixed lower threshold. Excluding the remaining terminal comparisons would give
\[
 \nuMin(p,e,2)=\left\lceil\frac{1+\log_p e+v_p(e)}2\right\rceil
\]
for every positive multiple $e$ of $p-1$. This is a proposed extension, not a proved formula in the remaining wild cases. For $r\ge1$, the next coefficient $k=(p^2+1)/2$ is prime to $p$ and remains outside the proved exact band. Theorem~\ref{thm:first-boundary-normal-form} forces any unequal comparison there to have $z_1=z_2=((p-1)/2)p^{r+3}$, opposite tangent scalars, and split residual representations with exchanged canonical lines. Corollary~\ref{cor:boundary-point-field-degrees} further forces canonical point-field degrees at least $p^{r+3}$ at level $p^{n-1}$ and $p^{r+5}$ at level $p^n$. The translation spaces persist along the actual point tower, so the last three relative steps all have degree $p^2$; when $r=1$, the successive total degrees are exactly $p^2,p^4,p^6$. Corollary~\ref{cor:boundary-nonsplit-cartan-image} further excludes containment of either actual finite image over the common tame good field in a principal nonsplit Cartan. Corollary~\ref{cor:boundary-nonsplit-self-twist} excludes the corresponding tame quadratic self-twist with irreducible intertwiner. Whether the remaining non-Cartan form occurs for actual full $p^{r+3}$-torsion remains open. The detailed reducible point-field profiles have their own explicitly stated local hypotheses and proofs; the full-class first-wild exclusion does not replace those local arguments. In particular the exchanged-high-line and large-wild-image restrictions remain useful as local comparison statements, while the graph valuation bound excludes the entire unequal-character comparison at $e=p(p-1)$.

For the universal level deformation at weights $(p^2+p+1,p^2)$, Proposition~\ref{prop:divisorial-full-cartan} gives the exact full principal Cartan invariants at every level: $e=f=p^{n-1}$ and residue field $k(t^{p^{n-1}})$. The separate scalar calculation at level four gives relative invariants $e=p^2,f=p$ between the scalar and full Cartan fixed fields. At $p=5,n=4$, the additional tame quotient of ramification index $12$ with unchanged residue field has absolute index exactly $6500$. Proposition~\ref{prop:cartan-torus-regularity} additionally proves that the actual quotient is regular, with regular Cartier reduced vertical divisor, at every nonzero point of the exceptional torus. Corollary~\ref{cor:cartan-torus-specialization} gives $78000\mid e_L$ for every local specialization of those level-four charts; if $e(L/K)\le12$, it follows that $e_K\ge6500$. These chart identities also obstruct an actual boundary comparison when its finite image is contained in a principal nonsplit Cartan: Corollary~\ref{cor:boundary-nonsplit-cartan-image} constructs an invariant-ring map and contradicts the required divisibility, including over arbitrary perfect residue fields. A general comparison need not have that image, and the divisorial value alone is not a bound for arbitrary normal models. At a normal surface point with a single vertical multiplicity $6500$, Lemma~\ref{lem:cartier-index-contact} shows that an index-$260$ specialization requires Cartier index divisible by $25$. The known cyclic quotient in Example~\ref{ex:noncartier-quotient-specialization} exhibits the corresponding denominator phenomenon at vertical multiplicity $1300$ and Cartier index $5$.

Under the actual-comparison and unit-preserving chart hypotheses of Proposition~\ref{prop:boundary-specialization-contact}, an upstairs specialization has contact divisible by $5$. A downstairs transverse curve may have precisely such contact. What remains is to analyze the non-Cartan images and formal models compatible with the actual finite-level isomorphism and unequal special characters. Proposition~\ref{prop:cartan-orbit-data} proves that the point-orbit sizes alone do not force Cartan containment; the full congruence image can be strictly larger and noncommutative. The exchanged residual lines do not themselves identify the two deformations as quadratic twists. The remaining degree-$p^6$ boundary configuration is neither constructed nor excluded.

For the first numerical case, $p=5$ and $e=260$, the index convention gives
\[
                         3\le\nuMin(5,260,2)\le4.
\]
Indeed, Theorem~\ref{thm:mixed-all-indices} supplies an unequal-character
$p^3$ comparison since $260>5^{6-3-1}=25$, whereas
Theorem~\ref{thm:all-index-upper} gives recovery from $p^5$-torsion
since $260\le5^4$. Recovery from actual $p^4$-torsion would give
index three; the existence of an unequal-character $p^4$ comparison
over an allowed field of index $260$ would give index four.
The nonsplit-Cartan exclusion restricts that undecided comparison
without deciding between these two values.

The tame all-prime and wild characteristic-three existence criteria allow the base field to vary. Corollary~\ref{cor:cyclotomic-full-p} settles the particular fields $\Q_p(\zeta_p)$, but these criteria do not classify which other fixed fields of a given index carry each example, nor do they determine all higher-height analogues. Proposition~\ref{prop:height-one-index} settles height one for every prime; the remaining higher-height questions have a different scope. Likewise, the character index concerns a group with potentially good graded pieces, whereas full rational torsion is a property of one actual representation. Keeping these quantifiers distinct is necessary when applying the theorems to a particular elliptic curve or Barsotti--Tate group.


The three necessary configurations produced by the primitive-grading reduction are all excluded by Corollary~\ref{pg-cor:exclusion}. The central criterion therefore closes that comparison argument as well as the shorter graph-valuation proof. The remaining index questions in this section concern absolute ramification divisible by $p$; the full-class prime-to-$p$ table has no unresolved interval.
